\bookchapter{What a Text File Is}{ch:01-what-a-text-file-is}

\begin{chapterintro}

Before you can process a text file, you need a clear picture of what one is. This chapter gives you that picture: bytes on a disk, an encoding that turns them into characters, and the newline that turns characters into lines.

\end{chapterintro}

\emph{The problem.} I open \texttt{notes.txt} in an editor and see three short lines. A hex viewer\index{hex viewer} shows the same file as a row of numbers.

\emph{The question.} Which one is the file, and why does a program that counts lines sometimes disagree with my eyes?

\section{Bytes First}\label{01-what-a-text-file-is:bytes-first}

A \textbf{file}\index{file} is a named sequence of \textbf{bytes}\index{bytes} stored on a disk. A byte is a number from 0 to 255; nothing more. The file system remembers the name and the length, and it doesn't care what the numbers mean.

A \textbf{text file}\index{text file} is a file whose bytes are meant to be read as characters. The rule that maps bytes to characters is the file's \textbf{encoding}\index{encoding}. Today that is almost always \textbf{UTF-8}\index{utf-8@UTF-8}, which uses one byte for plain English letters and two to four bytes for everything else.

\begin{callout}{note}

``Plain text'' is not a format of its own. It's an agreement between whoever wrote the bytes and whoever reads them: \emph{use UTF-8}. When the two sides disagree, you see garbage such as \texttt{café} instead of \texttt{café}.

\end{callout}

\codeneed{13}

Here is the whole idea in one picture. The same three lines, as the disk stores them and as an editor shows them:

\codeneed{10}\begin{codeblock}
\begin{Verbatim}[breaklines=false, fontsize=\fontsize{8.8}{8.68}\selectfont, fontfamily=diagrammono, formatcom=\diagramlines]
┌──────────────────────── notes.txt (31 bytes on disk) ────────────────────────┐
│ 74 68 65 20 63 61 74 20 73 61 74 0a 6f 6e 20 74 68 65 20 6d 61 74 0a 74 68 … │
└──────────────────────────────────────┬───────────────────────────────────────┘
                                       │  decode as UTF-8
                                       ▼
                      ┌─────────────────────────────────┐
                      │ the cat sat\n                   │
                      │ on the mat\n                    │
                      │ the end\n                       │
                      └─────────────────────────────────┘
\end{Verbatim}
\end{codeblock}

\section{Lines Are a Convention Too}\label{01-what-a-text-file-is:lines-are-a-convention-too}

A \textbf{line}\index{line} is the text between two \textbf{newline}\index{newline} characters.\index{line ending} The newline is an ordinary byte (\texttt{0a}, written \texttt{\textbackslash{}n} in code). Editors hide it and start a new row instead; to a program it's just another character.

\begin{callout}{tip}

When a line count looks one short, check whether the file ends with a newline. \texttt{wc\ -l} counts newline characters, not rows, so a last line without one isn't counted.

\end{callout}

\Needspace{13\baselineskip}

Different systems historically chose different line endings. You'll meet all three:

\begin{booktable}
\begin{longtable}{>{\raggedright\arraybackslash}p{\dimexpr 0.2650\linewidth-1.4699\tabcolsep\relax}>{\raggedright\arraybackslash}p{\dimexpr 0.1242\linewidth-1.7515\tabcolsep\relax}>{\raggedright\arraybackslash}p{\dimexpr 0.0807\linewidth-1.8386\tabcolsep\relax}>{\raggedright\arraybackslash}p{\dimexpr 0.5301\linewidth-0.9399\tabcolsep\relax}}
\tabletop
\rowcolor{tablehead}\thead{System} & \thead{Line ending} & \thead{Bytes} & \thead{Where you still see it} \\
\tablemid
\endfirsthead
\tabletop
\rowcolor{tablehead}\thead{System} & \thead{Line ending} & \thead{Bytes} & \thead{Where you still see it} \\
\tablemid
\endhead
\tablebottom
\endlastfoot
Unix, Linux, macOS & LF\index{line ending!LF} & \texttt{0a} & Almost every file on a server, every Git repository by default \\
Windows & CR LF\index{line ending!CR LF} & \texttt{0d\ 0a} & Files created by Notepad and older Windows tools; some CSV exports from
spreadsheets \\
Classic Mac OS (before 2001) & CR & \texttt{0d} & Very old files only; treat it as a curiosity \\
\end{longtable}
\end{booktable}

\begin{callout}{warning}

Never ``fix'' line endings on a file you don't own by saving it in another editor. A shell script with \texttt{CR\ LF} endings fails on Linux with a confusing \texttt{bad\ interpreter} error, and a changed file can silently break a checksum or a signature.

\end{callout}

\section{Lab: See the Bytes for Yourself}\label{01-what-a-text-file-is:lab-see-the-bytes-for-yourself}

\texttt{xxd} prints any file's bytes in hexadecimal, next to their characters.

\begin{enumerate}
\def\labelenumi{\arabic{enumi}.}
\item
  Run \texttt{xxd\ notes.txt}. \textbf{Expected:}

  \codeneed{2}\begin{codeblock}
  \begin{Verbatim}[formatcom=\outputstyle]
  00000000: 7468 6520 6361 7420 7361 740a 6f6e 2074  the cat sat.on t
  00000010: 6865 206d 6174 0a74 6865 2065 6e64 0a    he mat.the end.
  \end{Verbatim}
  \end{codeblock}
\item
  Count the \texttt{0a} bytes. \textbf{Expected:} three. Each \texttt{0a} is a newline, which \texttt{xxd} prints as a dot on the right, so the file has three lines.
\end{enumerate}

\section{Summary, Key Terms, and Review Questions}\label{01-what-a-text-file-is:summary-key-terms-and-review-questions}

\subsection*{Summary}\label{01-what-a-text-file-is:summary}

\begin{itemize}
\tightlist
\item
  A file is a named sequence of bytes; the disk stores numbers, not letters.
\item
  An \textbf{encoding} maps bytes to characters. Use \textbf{UTF-8} unless you have a reason not to.
\item
  A line ends at a \textbf{newline} byte. Line endings differ between systems.
\end{itemize}

\Needspace{10\baselineskip}

\subsection*{Key Terms}\label{01-what-a-text-file-is:key-terms}

\begin{booktable}
\begin{longtable}{>{\raggedright\arraybackslash}p{\dimexpr 0.1667\linewidth-1.6667\tabcolsep\relax}>{\raggedright\arraybackslash}p{\dimexpr 0.8333\linewidth-0.3333\tabcolsep\relax}}
\tabletop
\rowcolor{tablehead}\thead{Term} & \thead{Meaning} \\
\tablemid
\endfirsthead
\tabletop
\rowcolor{tablehead}\thead{Term} & \thead{Meaning} \\
\tablemid
\endhead
\tablebottom
\endlastfoot
File & A named sequence of bytes stored on a disk \\
Encoding & The rule that turns bytes into characters, such as UTF-8 \\
Newline & The character (\texttt{\textbackslash{}n}, byte \texttt{0a}) that ends a
line \\
\end{longtable}
\end{booktable}

\subsection*{Review Questions}\label{01-what-a-text-file-is:review-questions}

\begin{enumerate}
\def\labelenumi{\arabic{enumi}.}
\tightlist
\item
  Why can a text editor and a hex viewer show the ``same'' file so differently?
\item
  What does \texttt{wc\ -l} actually count?
\item
  What goes wrong when a writer and a reader disagree about the encoding?
\end{enumerate}

\begin{understandbox}

\textbf{You understand this when} you can explain why \texttt{café} takes 4 characters but 5 bytes in UTF-8.

\end{understandbox}
